\documentclass[11pt]{article}
\usepackage[margin=1in]{geometry}
\usepackage{amsmath,amssymb,amsthm,mathtools}
\usepackage{microtype}
\usepackage{url}

\newtheorem{theorem}{Theorem}
\newtheorem{lemma}[theorem]{Lemma}
\newtheorem{corollary}[theorem]{Corollary}
\theoremstyle{remark}
\newtheorem{remark}[theorem]{Remark}

\newcommand{\tS}{\widetilde{S}}
\newcommand{\Z}{\mathbb{Z}}
\newcommand{\Q}{\mathbb{Q}}
\newcommand{\SL}{\operatorname{SL}}

\title{A unipotent residue obstruction\\ for higher-weight generalized Dedekind sums}
\author{}
\date{15 November 2026}

\begin{document}
\maketitle

\begin{abstract}
At weight \(2\), a newform Dedekind sum is a homomorphism, so its image is a group.
At higher weight the same construction yields a period-polynomial cocycle, and evaluation of that cocycle need not be a homomorphism.
We prove a coefficient congruence, stable under a two-dimensional unipotent action modulo an odd prime, which forces a normalized value to miss a residue class.
On \(\Gamma_1(9)\), \(\Gamma_1(12)\), and \(\Gamma_1(15)\) the congruence applies to \((\chi_3,\chi_3,4)\), \((\chi_3,\chi_4,4)\), and \((\chi_3,\chi_5,7)\).
In each case the set of values is a proper subset of the additive group it generates, so it is not a group.
The second bullet of Tranbarger's Conjecture~1.15 therefore fails for these three triples.
The largest constants of containment printed for the first two triples remain valid.
The constant \(4\) printed for \((\chi_3,\chi_5,7)\) is already false on the finite set \(G_{50}(15)\) used to print the table.
One further triple on the same group, \((\chi_3,\chi_3,6)\), has the same shape: \(\tfrac{3}{10}\tS\) is an integer congruent to \(0\) or \(1\) modulo \(3\), the additive span is \(\tfrac{10}{3}\Z\), and the printed constant \(\tfrac{10}{3}\) remains the largest constant of containment.
A separate law modulo \(9\) sharpens the second example.
The exact image inside the surviving classes is not determined.
\end{abstract}

\section{The conjecture}
\label{sec:conj}

At weight \(2\), a newform Dedekind sum is a crossed homomorphism, so its image on \(\Gamma_1(q_1q_2)\) is a group \cite{Majure,SVY}.
At higher weight the same construction produces a period polynomial of degree \(n=k-2\) and only the cocycle \eqref{eq:cocycle} below.
Evaluation of that cocycle need not be a homomorphism.
For \((\chi_3,\chi_3,4)\), \((\chi_3,\chi_4,4)\), \((\chi_3,\chi_5,7)\), and \((\chi_3,\chi_3,6)\), one rescaling is an integer congruent to a square residue modulo \(3\), hence to \(0\) or \(1\).
Each image is then a proper subset of the additive group it generates.

The definitions and the conjecture are those of Tranbarger \cite{Tran}.
Numbering follows that text.
Let \(\chi_3\) be the primitive character modulo \(3\), with values \((0,1,-1)\) on \((0,1,2)\), and let \(\chi_4\) be the primitive character modulo \(4\), vanishing on even arguments, with \(\chi_4(1)=1\) and \(\chi_4(3)=-1\).
Both are odd.
The Legendre symbol \(\chi_5\) modulo \(5\) is even.
The parity condition \(\chi_1\chi_2(-1)=(-1)^k\) therefore allows \((\chi_3,\chi_3,4)\), \((\chi_3,\chi_4,4)\), and \((\chi_3,\chi_5,7)\), and it excludes \((\chi_3,\chi_3,3)\) and \((\chi_3,\chi_4,3)\).
These characters are real and equal to their own inverses.

For \(c>0\), Theorem~1.7 of the source is the finite sum
\[
S_{\chi_1,\chi_2,k}\begin{pmatrix}a&b\\c&d\end{pmatrix}
=\sum_{j\bmod c}\sum_{n\bmod q_1}
\overline{\chi_1}(n)\,\overline{\chi_2}(j)\,
B_1\Bigl(\frac{j}{c}\Bigr)
B_{k-1}\Bigl(\frac{aj}{c}+\frac{n}{q_1}\Bigr),
\]
with \(B_m\) the periodic Bernoulli function of Definition~1.1.
Definition~1.14 sets \(\tS(\gamma)=c^{k-2}S(\gamma)\).
The sum depends on \(\gamma\) only through \(a\) and \(c\).
For \(c=0\) one has \(\tS=0\) when \(k>2\).
For \(k=4\) one has \(\tS(-\gamma)=\tS(\gamma)\); for \(k=7\) one has \(\tS(-\gamma)=-\tS(\gamma)\), whenever both values are read off the same cusp.

On \(\Gamma_1(N)\) the nebentypus is trivial.
Lemma~5.1 therefore supplies, for each \(\gamma\), a polynomial \(h_\gamma\) of degree at most \(n=k-2\).
Lemma~5.2 is the right cocycle
\begin{equation}
\label{eq:cocycle}
h_{\gamma_1\gamma_2}=h_{\gamma_1}\big|_{2-k}\gamma_2+h_{\gamma_2},
\end{equation}
where \((f|_{-n}\gamma)(z)=(cz+d)^n f(\gamma z)\).
In particular \(h_{g^{-1}}=-(h_g|_{-n}g^{-1})\).
Formula~(12) evaluates the cocycle on an inverse.
The reading used below, which agrees with the finite sum for every matrix cited in this note, is the following.
If \(h_M(z)=\sum_{i=0}^{n}a_i(M)z^{n-i}\) and \(M=\begin{pmatrix}a&b\\c&d\end{pmatrix}\), then
\begin{equation}
\label{eq:eval}
\tS(M^{-1})=\sum_{i=0}^{n}a_i(M)\,d^{n-i}(-c)^{i}.
\end{equation}
The power is \((-c)^{k-2}\), with \(c\) the bottom-left entry of \(M\), not of \(M^{-1}\).
When \(k-2\) is even this equals \(c^{k-2}h_M(d/(-c))\).
When \(k-2\) is odd the two differ by a sign.
The matrix \(M=\begin{pmatrix}1&0\\-15&1\end{pmatrix}\) has \(c^5 h_M(d/(-c))=17000/3\), while the finite sum gives \(\tS(M^{-1})=\tS\begin{pmatrix}1&0\\15&1\end{pmatrix}=-17000/3\).

The source defines, for a positive integer \(j\),
\[
G_j(N)=\bigl\{\begin{pmatrix}a&b\\c&d\end{pmatrix}\in\Gamma_1(N):1\le a<jN,\ N\le c<jN\bigr\}.
\]
The tables of Section~6.1 record the largest rational \(r\ge 0\) such that \(\tS(G_{50}(q_1q_2))\subseteq r\Z\).
Conjecture~1.15, in the arXiv text, reads as follows.

\begin{quote}
We conjecture the following points regarding the image \(\tS_{\chi_1,\chi_2,k}(\Gamma_1(q_1q_2))\) for pairs of quadratic characters \(\chi_1\) and \(\chi_2\).
\begin{itemize}
\item The tables in Section~6.1 extend from \(G_{50}(q_1q_2)\) to \(\Gamma_1(q_1q_2)\) with equality.
\item For some \(d\) dividing \(2k-2\), either \(\tS=d\Z\) or \(\tS=d\Z/q_1\).
\end{itemize}
\end{quote}

The second bullet names only groups.
Theorem~\ref{thm:three} shows that each of the three images is a proper subset of its additive span, so none of them is a group, and the second bullet fails for these triples.
The first bullet is the sentence quoted above.
It is not the statement proved here.
For the first two triples the printed constant remains the largest positive rational \(r\) with the image contained in \(r\Z\).
For \((\chi_3,\chi_5,7)\) the printed constant \(4\) is already false on \(G_{50}(15)\), by the value in Theorem~\ref{thm:three}(iii); the largest such constant on the whole group is \(\tfrac43\).

Throughout, \(\langle I\rangle\) denotes the additive group generated by a set \(I\subseteq\Q\).
The index formula \([\SL_2(\Z):\Gamma_1(N)]=N^2\prod_{p\mid N}(1-1/p^2)\) gives the three indices \(72\), \(96\), and \(192\) used below.
The symmetric generating set of \(\SL_2(\Z)\) is \(\{S,T,S^{-1},T^{-1}\}\), with \(S=\begin{pmatrix}0&-1\\1&0\end{pmatrix}\) and \(T=\begin{pmatrix}1&1\\0&1\end{pmatrix}\).
A rewriting that omits the inverses does not present the group.
Generation is the Schreier lemma: once every Schreier generator is exhibited as a word in a proposed set and its inverses, that set generates.
The words are finite and are multiplied in the accompanying programs.
Those multiplications are integer arithmetic.
They are not a substitute for the cocycle induction.

\section{A unipotent congruence}
\label{sec:lemma}

\begin{lemma}
\label{lem:uni}
Let \(\ell\) be an odd prime, let \(q\) be a positive integer divisible by \(\ell\), let \(n\ge 1\), and let \(\lambda\in\Q^\times\).
Let \(\Gamma\le\SL_2(\Z)\) be a subgroup whose elements \(\begin{pmatrix}a&b\\c&d\end{pmatrix}\) satisfy \(q\mid c\) and \(a\equiv d\equiv 1\pmod{\ell}\).
Let \(h\) be a right cocycle of degree at most \(n\) on \(\Gamma\), in the sense of \eqref{eq:cocycle}, and write \(h_M(z)=\sum_{i=0}^{n}a_i(M)z^{n-i}\).
Set
\[
b_i(M)=\frac{q^{i+1}a_i(M)}{\lambda},\qquad r(M)=-\frac{c(M)}{q}.
\]
Suppose that on a generating set of \(\Gamma\) every \(b_i\) is an integer and, for a fixed \(\alpha\in\mathbb{F}_\ell\),
\begin{equation}
\label{eq:law}
b_0\equiv\alpha r^2,\qquad b_1\equiv -2\alpha r,\qquad b_i\equiv 0\quad(i\ge 2)\pmod{\ell}.
\end{equation}
Then \eqref{eq:law} holds for every element of \(\Gamma\), and every \(b_i\) remains an integer.
If the values attached to \(h\) satisfy \eqref{eq:eval}, the normalized quantity
\[
U(M^{-1})=\frac{q}{\lambda}\,\tS(M^{-1})
\]
is an integer, and
\begin{equation}
\label{eq:U}
U(M^{-1})\equiv -\alpha\, r(M)^2\pmod{\ell}.
\end{equation}
Since inversion permutes \(\Gamma\), the same congruence holds for \(U(\gamma)\) with \(r(M)\) replaced by \(c(\gamma)/q\).
If \(\alpha\not\equiv 0\), the residues of \(U\) therefore lie among the \((\ell+1)/2\) values of \(-\alpha\) times a square in \(\mathbb{F}_\ell\), and at least \((\ell-1)/2\) residue classes do not occur.
\end{lemma}

The additive group generated by the values of \(U\) is not part of the lemma.
Integrality of the coefficients \(b_i\) on the generating set is an input, not a conclusion about arbitrary polynomials.
The lemma does not require \(\ell\) to be the full conductor, only that \(\ell\) divide \(q\) and that the diagonal entries be \(1\) modulo \(\ell\).

\begin{proof}
Rescale the polynomial by
\[
F_M(X)=\frac{q^{n+1}}{\lambda}\,h_M\Bigl(\frac{X}{q}\Bigr)=\sum_{i=0}^{n}b_i(M)X^{n-i}.
\]
For \(g=\begin{pmatrix}a&b\\c&d\end{pmatrix}\) write \(g_q=\begin{pmatrix}a&qb\\c/q&d\end{pmatrix}\).
Then \(\det g_q=1\).
The hypothesis \(q\mid c\) makes \(g_q\) an integral matrix.
The weight \(-n\) slash intertwines with this conjugation.
Indeed
\[
\frac{q^{n+1}}{\lambda}\,(h|_{-n}g)\Bigl(\frac{X}{q}\Bigr)
=\frac{q}{\lambda}\,(cX+qd)^n\,
h\Bigl(\frac{aX+qb}{cX+qd}\Bigr),
\]
while \(g_q\cdot X=(aX+qb)/((c/q)X+d)\) gives the same expression for \((F_M|_{-n}g_q)(X)\).
Explicitly,
\begin{equation}
\label{eq:Faction}
(F_M|_{-n}g_q)(X)
=\sum_{i=0}^{n}b_i(M)\,(aX+qb)^{n-i}\Bigl(\frac{c}{q}X+d\Bigr)^{i}.
\end{equation}
If the \(b_i(M)\) are integers, every coefficient on the right is an integer.
Addition and negation preserve \(\Z[X]\).
The cocycle becomes \(F_{Mg}=(F_M|_{-n}g_q)+F_g\), and \(F_{g^{-1}}=-(F_g|_{-n}g_q^{-1})\).
The inverse matrix \(g_q^{-1}\) is integral of determinant \(1\).
Thus integrality on a generating set passes to every word and every inverse, hence to \(\Gamma\).

Now reduce modulo \(\ell\).
The conditions \(a\equiv d\equiv 1\pmod{\ell}\), \(\ell\mid q\), and \(r(g)=-c(g)/q\) give
\[
g_q\equiv\begin{pmatrix}1&0\\-r(g)&1\end{pmatrix}\pmod{\ell}.
\]
In \eqref{eq:Faction}, every term with \(i\ge 2\) is multiplied by \(b_i(M)\).
On the generating set, and then by the inductive argument below on the whole group, those coefficients are \(0\) modulo \(\ell\), so they contribute nothing modulo \(\ell\) to any output coefficient.
The remaining terms are
\[
b_0(aX+qb)^n+b_1(aX+qb)^{n-1}\Bigl(\frac{c}{q}X+d\Bigr)
\equiv b_0 X^n+b_1 X^{n-1}(-r(g)X+1)\pmod{\ell}.
\]
The leading pair therefore moves by
\begin{equation}
\label{eq:unipotent}
(b_0,b_1)\longmapsto(b_0-r(g)\,b_1,\ b_1),
\end{equation}
and the coefficients of \(X^{n-i}\) for \(i\ge 2\) stay \(0\) modulo \(\ell\).
The hypothesis that \(\ell\) is odd keeps the linear condition nontrivial: if \(\ell=2\), then \(-2\alpha\equiv 0\) and \eqref{eq:law} would force \(b_1\equiv 0\).

Let \(M\) and \(g\) both satisfy \eqref{eq:law}, and write \(r_M=r(M)\), \(r_g=r(g)\).
Adding \(F_g\) to the image under \eqref{eq:unipotent} produces
\begin{align*}
b_0(Mg)
&\equiv\alpha r_M^2-r_g(-2\alpha r_M)+\alpha r_g^2
=\alpha(r_M+r_g)^2,\\
b_1(Mg)
&\equiv -2\alpha r_M-2\alpha r_g
=-2\alpha(r_M+r_g)\pmod{\ell}.
\end{align*}
Matrix multiplication gives \(c(Mg)=c(M)a(g)+d(M)c(g)\), so
\[
r(Mg)=r_M\,a(g)+d(M)\,r_g\equiv r_M+r_g\pmod{\ell},
\]
where the last step uses \(a(g)\equiv d(M)\equiv 1\pmod{\ell}\).
Thus \(Mg\) satisfies \eqref{eq:law}.
The identity has \(r=0\) and \(F=0\), so it satisfies the law.

For the inverse, \(r(g^{-1})=-r_g\) and \(g_q^{-1}\equiv\begin{pmatrix}1&0\\r_g&1\end{pmatrix}\).
The ratio of the acting matrix \(g^{-1}\) is \(-r_g\), so \eqref{eq:unipotent} is applied with that ratio:
\[
\bigl(\alpha r_g^2,\ -2\alpha r_g\bigr)
\longmapsto
\bigl(\alpha r_g^2+r_g(-2\alpha r_g),\ -2\alpha r_g\bigr)
=\bigl(-\alpha r_g^2,\ -2\alpha r_g\bigr).
\]
Negating yields \((\alpha r_g^2,\ 2\alpha r_g)\), which is \eqref{eq:law} at ratio \(-r_g\).
Higher coefficients remain \(0\) modulo \(\ell\).
By induction on word length, \eqref{eq:law} holds on \(\Gamma\).

Finally let \(r=r(M)\).
The relation \(r=-c/q\) is \(-\,c=qr\), so \((-c)^i/q^i=r^i\).
Substituting \(a_i=\lambda b_i/q^{i+1}\) into \eqref{eq:eval} gives
\[
U(M^{-1})=\sum_{i=0}^{n}b_i(M)\,d^{n-i}r^{i}.
\]
Modulo \(\ell\) one has \(d\equiv 1\) and \(b_i\equiv 0\) for \(i\ge 2\), hence
\[
U(M^{-1})\equiv b_0+b_1 r\equiv\alpha r^2+(-2\alpha r)r=-\alpha r^2\pmod{\ell}.
\]
The map \(r\mapsto r^2\) on \(\mathbb{F}_\ell\) takes \((\ell+1)/2\) values.
If \(\alpha\not\equiv 0\), multiplication by \(-\alpha\) does not enlarge that set.
\end{proof}

The three applications in the next section all take \(\ell=3\).
The pairs \((n,q,\lambda,\alpha)\) and the resulting normalizations are

\begin{center}
\begin{tabular}{lccccc}
\hline
triple & \(n\) & \(q\) & \(\lambda\) & \(\alpha\) & \(U\) \\
\hline
\((\chi_3,\chi_3,4)\) & \(2\) & \(9\) & \(18\) & \(1\) & \(\tS/2\) \\
\((\chi_3,\chi_4,4)\) & \(2\) & \(3\) & \(2\) & \(2\) & \(3\tS/2\) \\
\((\chi_3,\chi_5,7)\) & \(5\) & \(3\) & \(4\) & \(2\) & \(3\tS/4\) \\
\hline
\end{tabular}
\end{center}

In the last two rows, \(-\alpha=-2\equiv 1\pmod{3}\), so \eqref{eq:U} says \(U\equiv r^2\pmod{3}\).
In the first row it says \(U\equiv -r^2\pmod{3}\).
Combined with integrality of \(U\), that is the congruence modulo \(6\) stated below.
None of the three span computations is implied by Lemma~\ref{lem:uni}.

\section{Three images}
\label{sec:three}

\begin{theorem}
\label{thm:three}
Let \(I(\chi_1,\chi_2,k)\) denote the set of values of \(\tS_{\chi_1,\chi_2,k}\) on \(\Gamma_1(q_1q_2)\).
\begin{enumerate}
\item The set \(I(\chi_3,\chi_3,4)\) is contained in \(\{n\in\Z:n\equiv 0\text{ or }4\pmod{6}\}\), contains \(-2\), and does not contain \(2\).
Its additive span is \(2\Z\).
The largest positive rational \(r\) with \(I(\chi_3,\chi_3,4)\subseteq r\Z\) is \(r=2\).
\item The set \(I(\chi_3,\chi_4,4)\) is contained in \(\tfrac23\Z\), and its additive span is exactly \(\tfrac23\Z\).
For every value, \(\tfrac32\tS\) is an integer congruent to \(0\) or \(1\) modulo \(3\), so \(\tfrac43\) and \(-\tfrac23\) do not occur.
The largest positive rational \(r\) with the image contained in \(r\Z\) is \(r=\tfrac23\).
\item The set \(I(\chi_3,\chi_5,7)\) is contained in \(\tfrac43\Z\), and its additive span is exactly \(\tfrac43\Z\).
For every value, \(\tfrac34\tS\) is an integer congruent to \(0\) or \(1\) modulo \(3\), so \(\tfrac83\) and \(-\tfrac43\) do not occur.
The largest positive rational \(r\) with the image contained in \(r\Z\) is \(r=\tfrac43\).
Separately,
\[
\tS_{\chi_3,\chi_5,7}\begin{pmatrix}1&-1\\15&-14\end{pmatrix}=-\frac{17000}{3},
\]
and this matrix lies in \(G_{50}(15)\).
The value does not lie in \(4\Z\).
\end{enumerate}
In all three cases the image is a proper subset of its additive span, hence is not a group.
Every set named in the second bullet of Conjecture~1.15 is a group.
That bullet fails for all three triples.
\end{theorem}

The positive divisors of \(2\cdot 4-2=6\) are \(1,2,3,6\), and \(q_1=3\), so the groups allowed by the second bullet in the weight-\(4\) cases are \(\Z\), \(\tfrac13\Z\), \(2\Z\), \(\tfrac23\Z\), \(3\Z\), and \(6\Z\).
The same list with the lattice \(\tfrac43\Z\) in place of the weight-\(4\) spans is not needed: once the image is not a group, it equals none of them.

\subsection{Level \(9\), weight \(4\)}
\label{sec:9}

For \(M=\begin{pmatrix}a&b\\9m&d\end{pmatrix}\in\Gamma_1(9)\) write
\[
h_M=Ax^2+\frac{B_1}{9}x+\frac{C_1}{9}
\]
whenever \((A,B_1,C_1)\) is an integral triple.
Formula \eqref{eq:eval} becomes
\[
\tS(M^{-1})=Ad^2-B_1\,d\,m+9C_1 m^2.
\]
The rescaling of Lemma~\ref{lem:uni} with \((q,\lambda)=(9,18)\) is
\[
(b_0,b_1,b_2)=\Bigl(\frac{A}{2},\frac{B_1}{2},\frac{9C_1}{2}\Bigr),\qquad r=-m.
\]
Thus \(U=\tS/2\), and integrality of \((b_0,b_1,b_2)\) is evenness of \((A,B_1,C_1)\).

The group \(\Gamma_1(9)\) is generated by the seven matrices \(\mathtt{Gamma1(9).gens()}\), stored with their period polynomials in the repository \cite{NFDS}.
The polynomials were compared with the cocycle definition of \(h_M\) at cusps of denominator \(9\), and \eqref{eq:eval} reproduces the finite sum on each inverse.
The coordinates are

\begin{center}
\begin{tabular}{ccc}
\hline
\(M\) & \((A,B_1,C_1)\) & \(m=c/9\) \\
\hline
\((1,1;0,1)\) & \((0,0,0)\) & \(0\) \\
\((-8,1;-9,1)\) & \((2,-4,0)\) & \(-1\) \\
\((19,-4;81,-17)\) & \((522,-1980,208)\) & \(9\) \\
\((10,-3;27,-8)\) & \((-174,876,-124)\) & \(3\) \\
\((-26,11;-45,19)\) & \((-220,1672,-352)\) & \(-5\) \\
\((37,-16;81,-35)\) & \((-612,4752,-1024)\) & \(9\) \\
\((19,-12;27,-17)\) & \((156,-1716,526)\) & \(3\) \\
\hline
\end{tabular}
\end{center}

Every entry is an even integer.
Reducing modulo \(3\), one finds \(A\equiv 2m^2\) and \(B_1\equiv -2m\).
Dividing by \(2\) modulo \(3\), which is multiplication by \(2\), rewrites these as \(b_0\equiv r^2\) and \(b_1\equiv -2r\), with the third coordinate \(0\) modulo \(3\).
This is \eqref{eq:law} with \(\alpha=1\).
The check is finite: for the second row, \(2-2(-1)^2=0\) and \(-4+2(-1)=-6\).
The other rows are the same arithmetic.

The index \([\SL_2(\Z):\Gamma_1(9)]\) is \(72\).
Enumerating right cosets by bottom rows modulo \(9\) produces \(19\) Schreier generators relative to \(\{S,T,S^{-1},T^{-1}\}\), one of which is the identity.
Each of the other eighteen equals a word of length at most \(5\) in the seven matrices and their inverses.
The words are the file \texttt{certificates/gamma1\_9\_words.txt}, and the products are recomputed by \texttt{src/check\_words.py}.
The Schreier lemma gives generation.
Lemma~\ref{lem:uni} therefore applies to all of \(\Gamma_1(9)\).

The conclusion \(U\equiv -r^2\pmod{3}\) means \(\tS/2\equiv -(c/9)^2\pmod{3}\) after the square is moved to the matrix being evaluated.
Since \(U\) is an integer, \(\tS=2U\) is an even integer and
\[
\tS(\gamma)\equiv -2\Bigl(\frac{c}{9}\Bigr)^2\pmod{6}.
\]
Squares modulo \(3\) are \(0\) and \(1\).
The class \(2\pmod{6}\) would require \(U\equiv 1\pmod{3}\), which is the class omitted by \(-r^2\).
No value is congruent to \(2\) modulo \(6\).
The identity gives \(0\).
The matrix \(\begin{pmatrix}1&0\\9&1\end{pmatrix}\) lies in \(G_{50}(9)\) and has \(\tS=-2\), by the case \(m=1\) of Theorem~\ref{thm:family} in the appendix, or by the finite sum at \(c=9\).
Thus \(-2\) lies in the image, the image lies in \(2\Z\), and the span is exactly \(2\Z\).
The integer \(2\) is excluded, so the image is not \(2\Z\).
No positive rational larger than \(2\) contains \(-2\) in its multiples.

\subsection{Level \(12\), weight \(4\)}
\label{sec:12}

For \(M=\begin{pmatrix}a&b\\c&d\end{pmatrix}\in\Gamma_1(12)\) write \(c=3\,\mathrm{cpp}\) and
\[
h_M=\frac{2A}{3}x^2+\frac{2B}{9}x+\frac{2C}{27},
\]
so that \((A,B,C)=\bigl(\tfrac32 c_2,\tfrac92 c_1,\tfrac{27}{2}c_0\bigr)\) when \(h_M=c_2x^2+c_1x+c_0\).
On \(\Gamma_1(12)\) the integer \(\mathrm{cpp}\) is divisible by \(4\).
Formula \eqref{eq:eval} gives
\[
\tS(M^{-1})=\frac23\bigl(Ad^2-B\,d\,\mathrm{cpp}+C\,\mathrm{cpp}^2\bigr).
\]
The rescaling with \((q,\lambda)=(3,2)\) is \((b_0,b_1,b_2)=(A,B,C)\) and \(r=-\mathrm{cpp}\), so \(U=\tfrac32\tS\).
The slash in these coordinates, by a matrix with bottom-left entry \(3\,\mathrm{cpp}\), is
\begin{align*}
A'&=Aa^2+Ba\,\mathrm{cpp}+C\,\mathrm{cpp}^2,\\
B'&=6Aab+B(ad+bc)+2C\,\mathrm{cpp}\,d,\\
C'&=9Ab^2+3Bbd+Cd^2.
\end{align*}
Integral coordinates remain integral.
This is the integrality input of the lemma, seen in coordinates.

The nine matrices \(\mathtt{Gamma1(12).gens()}\) have coordinates

\begin{center}
\begin{tabular}{ccc}
\hline
\(M\) & \((A,B,C)\) & \(\mathrm{cpp}\) \\
\hline
\((1,1;0,1)\) & \((0,0,0)\) & \(0\) \\
\((-11,1;-12,1)\) & \((8,-4,0)\) & \(-4\) \\
\((49,-9;60,-11)\) & \((560,-616,168)\) & \(20\) \\
\((13,-3;48,-11)\) & \((-1087,1426,-471)\) & \(16\) \\
\((13,-4;36,-11)\) & \((-72,120,-48)\) & \(12\) \\
\((61,-25;144,-59)\) & \((-1800,4422,-2715)\) & \(48\) \\
\((13,-6;24,-11)\) & \((-40,110,-75)\) & \(8\) \\
\((25,-16;36,-23)\) & \((-216,840,-816)\) & \(12\) \\
\((37,-27;48,-35)\) & \((1127,-4862,5247)\) & \(16\) \\
\hline
\end{tabular}
\end{center}

Each non-identity polynomial agrees with the finite sum at the cusps \(1/12\), \(5/12\), \(7/12\), \(11/12\) and at least one of \(5/24\), \(7/24\).
Formula \eqref{eq:eval} reproduces the finite sum on each inverse.
The nine triples satisfy \(C\equiv 0\), \(B\equiv\mathrm{cpp}\), and \(A\equiv 2\,\mathrm{cpp}^2\pmod{3}\).
With \(r=-\mathrm{cpp}\) and \(\alpha=2\), one has \(\alpha r^2=2\,\mathrm{cpp}^2\) and \(-2\alpha r=4\,\mathrm{cpp}\equiv\mathrm{cpp}\pmod{3}\).
Thus \(A\equiv\alpha r^2\) and \(B\equiv -2\alpha r\).
This is \eqref{eq:law}.
For the second row, \(8-2\cdot 16=-24\), \(-4-(-4)=0\), and \(C=0\).

The index is \(96\).
The same symmetric set produces \(27\) Schreier generators.
The \(26\) nontrivial ones are words of length at most \(5\) in the nine matrices and their inverses, listed in \texttt{certificates/gamma1\_12\_words.txt} and multiplied back by \texttt{src/check\_words.py}.
Lemma~\ref{lem:uni} gives \(U\equiv -2r^2\equiv \mathrm{cpp}^2\pmod{3}\), hence \(U\equiv 0\) or \(1\pmod{3}\), never \(2\).

The finite sum returns
\[
\tS\begin{pmatrix}1&-1\\12&-11\end{pmatrix}=-\frac{16}{3},
\qquad
\tS\begin{pmatrix}-11&3\\-48&13\end{pmatrix}=-\frac{2254}{3}.
\]
The first matrix is the inverse of the second generator, and \eqref{eq:eval} agrees.
The second is the inverse of \((13,-3;48,-11)\), and \eqref{eq:eval} agrees again.
The greatest common divisor of these two rationals is \(\tfrac23\).
Combined with the containment \(I\subseteq\tfrac23\Z\) coming from integrality of \(U\), the span is \(\tfrac23\Z\).
Now \(\tfrac32\cdot\tfrac43=2\equiv 2\pmod{3}\) and \(\tfrac32\cdot(-\tfrac23)=-1\equiv 2\pmod{3}\), so neither rational lies in the image.
No positive rational larger than \(\tfrac23\) has the same containment.

\subsection{Level \(15\), weight \(7\)}
\label{sec:15}

The matrix \(\begin{pmatrix}1&-1\\15&-14\end{pmatrix}\) lies in \(\Gamma_1(15)\) and in \(G_{50}(15)\): the determinant is \(1\), the diagonal entries are \(1\) modulo \(15\), \(1\le 1<750\), and \(15\le 15<750\).
The finite sum of Theorem~1.7, with the degree-\(6\) periodic Bernoulli function and with \(\chi_5\) the Legendre symbol, equals \(-17000/3\).
This rational does not lie in \(4\Z\).
The cell \(4\) in the Section~6.1 table is therefore already false on the set from which the table was computed.
The rest of this subsection identifies the actual span and the residue obstruction.

Bottom rows modulo \(15\) with \(\gcd(c,d,15)=1\) number \(192\), the index \([\SL_2(\Z):\Gamma_1(15)]\).
Right cosets enumerated along \(\{S,T,S^{-1},T^{-1}\}\) give \(47\) Schreier generators, all in \(\Gamma_1(15)\), one of them the identity.
They generate the group.
On each of the \(46\) non-identity generators, Lemma~5.1 supplies a polynomial of degree at most \(5\).
The polynomial is the interpolation of
\[
h_\gamma(\alpha)=\widehat{S}(\alpha)-(c\alpha+d)^5\widehat{S}(\gamma\alpha)
\]
through six cusps whose denominators are divisible by \(15\), checked at a seventh cusp.
In the sums that produce these values, no denominator exceeds \(4515\).
Write \(h=\sum_{n=0}^{5}a_n x^{5-n}\) and \(b_n=3^{n+1}a_n/4\).
Each of the \(46\) polynomials has every \(b_n\) integral.
Equivalently, each lies in the space \(\mathcal{P}(7;4,3)\) of Definition~1.16.
One polynomial is short enough to inspect directly.
The Schreier matrix \(\begin{pmatrix}1&0\\-15&1\end{pmatrix}\) has
\[
(b_0,\ldots,b_5)=(4250,-3200,900,-120,0,0),
\]
so
\[
h=\frac{17000}{3}x^5-\frac{12800}{9}x^4+\frac{400}{3}x^3-\frac{160}{27}x^2.
\]
The four nonzero products \(3^{n+1}a_n\) are \(17000\), \(-12800\), \(3600\), and \(-480\), each divisible by \(4\).
The full list is the lines beginning with \texttt{P} in \texttt{certificates/verify-k7-log.txt}.
There are \(46\) such lines.
The identity contributes the zero polynomial.

This integrality is the input \(b_i\in\Z\) of Lemma~\ref{lem:uni}, with \((q,\lambda)=(3,4)\) and \(n=5\).
Once it is known on the Schreier generators, the integral matrix \(g_q\) preserves \(\Z[X]\) by \eqref{eq:Faction}, so the cocycle propagates integrality to every word.
The general fact that all of \(\mathcal{P}(7;4,3)\) is closed under the slash is a stronger statement than this argument uses.
It is the content of Lemma~6.2 in the source; the displayed monomial degree in that lemma's expansion should be \(x^{k-2-j}\), while the divisibility itself is the one in Definition~1.16.
Evaluation through \eqref{eq:eval} then puts every value of \(\tS\) in \(\tfrac43\Z\).
The case \(c=0\) is the value \(0\).

The same \(46\) polynomials satisfy \eqref{eq:law} with \(\alpha=2\) and \(r=-c/3\).
For the specimen above, \(r=5\), and \(4250\equiv 2\equiv 2\cdot 5^2\), \(-3200\equiv 1\equiv 2\cdot 5\), while the four lower integers are \(0\) modulo \(3\).
The shape \(b_1\equiv 2r\) agrees with \(b_1\equiv -2\alpha r=-4r\) because \(-4\equiv 2\pmod{3}\).
The program \texttt{src/verify\_refinements.py} repeats the congruence on all \(46\) lines.
Lemma~\ref{lem:uni} gives
\[
U=\frac34\tS\equiv -2r^2\equiv r^2\pmod{3},
\]
so \(U\) is \(0\) or \(1\) modulo \(3\), never \(2\).
The identity gives the class \(0\).
The matrix \(\begin{pmatrix}1&0\\15&1\end{pmatrix}\) has \(\tS=-17000/3\), hence \(U=-4250\equiv 1\pmod{3}\).

The matrix \(\begin{pmatrix}16&-9\\105&-59\end{pmatrix}\) lies in \(\Gamma_1(15)\), and the finite sum returns \(-1759326764/3\).
The greatest common divisor of this rational and \(-17000/3\) is \(\tfrac43\).
The span contains \(\tfrac43\Z\).
With the containment, the span equals \(\tfrac43\Z\), and no larger positive rational has the same containment.
The class \(U\equiv 2\pmod{3}\) is empty, so \(\tfrac83=\tfrac43\cdot 2\) and \(-\tfrac43=\tfrac43\cdot(-1)\) lie in the span and not in the image.

\section{A square law modulo \(9\)}
\label{sec:mod9}

Lemma~\ref{lem:uni} sees only the prime \(3\).
The second example has a sharper congruence.
It is not a consequence of the lemma.

\begin{theorem}
\label{thm:mod9}
Keep the coordinates \((A,B,C)\) and \(\mathrm{cpp}=c/3\) of Section~\ref{sec:12}.
For every \(M=\begin{pmatrix}a&b\\3\,\mathrm{cpp}&d\end{pmatrix}\in\Gamma_1(12)\),
\begin{equation}
\label{eq:mod9}
A\equiv 5\,\mathrm{cpp}^2,\qquad B\equiv \mathrm{cpp}\,d,\qquad C\equiv d-1\pmod{9}.
\end{equation}
Consequently \(U(M^{-1})=\tfrac32\tS(M^{-1})\) satisfies \(U\equiv 4\,\mathrm{cpp}^2\pmod{9}\).
The right-hand side is a square modulo \(9\), so \(U\) is congruent to \(0\), \(1\), \(4\), or \(7\) modulo \(9\), and never to \(3\) or \(6\).
In particular the integers \(2\) and \(4\) do not lie in \(I(\chi_3,\chi_4,4)\).
Both are allowed by the congruence modulo \(3\).
All four square classes occur.
\end{theorem}

\begin{proof}
The coordinate law implies the value formula before any cocycle is used.
Reducing the definition of \(U\) gives
\[
U\equiv 5\,\mathrm{cpp}^2\,d^2-(\mathrm{cpp}\,d)\,d\,\mathrm{cpp}+(d-1)\mathrm{cpp}^2
=\mathrm{cpp}^2(4d^2+d-1)\pmod{9}.
\]
Every diagonal entry of a matrix in \(\Gamma_1(12)\) is \(1\) modulo \(12\), so \(d=1+12\beta\) and \(d\equiv 1+3\beta\pmod{9}\).
Then \(d^2\equiv 1+6\beta\pmod{9}\) and
\[
4d^2+d-1\equiv 4+6\beta+1+3\beta-1=4+9\beta\equiv 4\pmod{9}.
\]
Hence \(U\equiv 4\,\mathrm{cpp}^2=(2\,\mathrm{cpp})^2\pmod{9}\).
The squares modulo \(9\) are \(0,1,4,7\).

The nine matrices of Section~\ref{sec:12}, and the nine inverse polynomials \(h_{g^{-1}}=-(h_g|_{-2}g^{-1})\), satisfy \eqref{eq:mod9}.
For the second row, \(\mathrm{cpp}=-4\equiv 5\), \(d=1\), \(A=8\), \(B=-4\), and \(C=0\): one has \(5\cdot 25=125\equiv 8\), \(5\cdot 1\equiv -4\), and \(d-1\equiv 0\pmod{9}\).
The other rows, and the nine inverses, are the same finite reduction.
These eighteen matrices generate \(\Gamma_1(12)\) as a monoid.

The law is preserved by right multiplication of any two elements that satisfy it.
Let \(M\) and \(g=\begin{pmatrix}a&b\\3q&\delta\end{pmatrix}\) both satisfy \eqref{eq:mod9}, with \(a\equiv\delta\equiv 1\pmod{12}\) and \(\det=1\).
Write \(p=\mathrm{cpp}(M)\).
An error of \(9\) in \(A\), \(B\), or \(C\) is sent by the slash formulae to an error of \(9\) in each output coordinate: every extra term carries an explicit factor \(9\), \(6\cdot 9\), \(3\cdot 9\), or \(2\cdot 9\).
The output modulo \(9\) may therefore be computed from the representatives \(A=5p^2\), \(B=pd\), \(C=d-1\) and \(A_g=5q^2\), \(B_g=q\delta\), \(C_g=\delta-1\).
The new bottom data are \(p_{\mathrm{new}}=pa+dq\) and \(d_{\mathrm{new}}=3pb+d\delta\).

The leading coordinate of the product is
\[
A_{\mathrm{new}}=5p^2 a^2+pda\,q+(d-1)q^2+5q^2,
\]
and \(5p_{\mathrm{new}}^2=5(pa+dq)^2\), so
\begin{align*}
A_{\mathrm{new}}-5p_{\mathrm{new}}^2
&=pda\,q-10\,pda\,q+(d+4)q^2-5d^2 q^2\\
&=-9\,pda\,q+q^2(4+d-5d^2).
\end{align*}
The first summand is divisible by \(9\).
Substituting \(d=1+12\beta\) gives \(d\equiv 1+3\beta\) and \(d^2\equiv 1+6\beta\pmod{9}\), hence
\[
4+d-5d^2\equiv 4+(1+3\beta)-5(1+6\beta)=9\beta-27\beta\equiv 0\pmod{9}.
\]

The middle coordinate is
\[
B_{\mathrm{new}}=30p^2 ab+pd(a\delta+3bq)+2(d-1)q\delta+q\delta.
\]
Subtracting \(p_{\mathrm{new}}d_{\mathrm{new}}=(pa+dq)(3pb+d\delta)\) cancels \(pd\,a\delta\) against \(pa\,d\delta\) and cancels \(3pdbq\) against itself, and leaves
\[
B_{\mathrm{new}}-p_{\mathrm{new}}d_{\mathrm{new}}=27p^2 ab+q\delta(2d-1-d^2).
\]
The first summand is divisible by \(9\).
The same expansion \(d=1+12\beta\) gives \(2d-1-d^2=-144\beta^2\).

For the constant coordinate the term \(9Ab^2\) is \(0\) modulo \(9\).
With the representatives, the constant coordinate of the product is congruent to
\[
3pdb\,\delta+(d-1)\delta^2+(\delta-1)\pmod{9},
\]
and \(d_{\mathrm{new}}-1=3pb+d\delta-1\).
Subtracting and regrouping \((d-1)\delta^2-\,(d-1)\delta\) produces the congruence
\begin{equation}
\label{eq:C}
C_{\mathrm{new}}-(d_{\mathrm{new}}-1)\equiv 3pb(d\delta-1)+(d-1)\delta(\delta-1)\pmod{9}.
\end{equation}
Both \(d-1\) and \(\delta-1\) are divisible by \(12\), so \(d\delta-1\) is divisible by \(12\) and the first summand by \(36\), while \((d-1)(\delta-1)\) is divisible by \(144\).

Every word in the eighteen matrices therefore satisfies \eqref{eq:mod9}, and every value of \(U\) is a square modulo \(9\).
The four classes occur on
\begin{align*}
U(1,0;0,1)&=0,\\
U(-11,1;-12,1)&=-8\equiv 1,\\
U(13,-6;24,-11)&=40\equiv 4,\\
U(49,-9;60,-11)&=-560\equiv 7\pmod{9},
\end{align*}
all computed from \eqref{eq:eval}.
The residues \(3\) and \(6\) are \(0\) modulo \(3\), so Theorem~\ref{thm:three}(ii) did not exclude them.
Their exclusion removes \(\tS=2\) and \(\tS=4\).
It does not determine \(U\) inside the four square classes.
\end{proof}

The right-hand side of \eqref{eq:C} is a congruence because \(9Ab^2\) was discarded.
It is not an equality of integers.

\section{What remains open}
\label{sec:open}

The images in Theorem~\ref{thm:three} and Corollary~\ref{cor:k6} are not classified inside the classes that survive.
For \((\chi_3,\chi_3,4)\) the obstruction of Theorem~\ref{thm:three}(i) allows every integer congruent to \(0\) or \(4\) modulo \(6\), and the note does not decide which of those integers occur.
A coefficient quotient on the state \((A,B_1,C_1,m,d)\), reduced modulo \(16\) and modulo \(36\) and acted on by the fourteen matrices appearing in the level-\(9\) word file, meets every even residue modulo \(16\) and every residue modulo \(36\) congruent to \(0\) or \(4\) modulo \(6\).
The reachable counts are \(24576\) states modulo \(16\) and \(11664\) states modulo \(36\).
Forty exact words of length \(6\) match the reduced action.
This enumeration shows that those two moduli do not cut the image further.
It is not a description of the image.
The same limitation holds for \(U\) modulo \(9\) in Theorem~\ref{thm:mod9}, for \(\tfrac34\tS\) modulo \(3\) in Theorem~\ref{thm:three}(iii), and for \(\tfrac{3}{10}\tS\) modulo \(3\) in Corollary~\ref{cor:k6}: both allowed classes occur, and the class inside them is not determined.

No claim is made about the cells of the Section~6.1 tables other than the four treated in Theorem~\ref{thm:three} and Corollary~\ref{cor:k6}.
The equality \(\tS_{\chi_3,\chi_5,3}(\Gamma_1(15))=4\Z\) is not a theorem of this note.
A scan that sees no positive value below \(94\) in a bounded range of \(c\) is an experiment, not a minimal-value theorem.
Lemma~\ref{lem:uni} applies to a further triple only when some odd prime \(\ell\) divides the step \(q\) in \(r=-c/q\), the diagonal is \(1\) modulo \(\ell\), and the initial coefficients on a generating set obey \eqref{eq:law} for some \(\alpha\not\equiv 0\).
Those initial coefficients are computed here case by case.
No general formula for them, at arbitrary weight, is proved.

\section{Certificates and attribution}
\label{sec:attr}

The inductions in Lemma~\ref{lem:uni} and Theorem~\ref{thm:mod9} are the arguments above.
The finite inputs are exact rational or integer arithmetic in the Python standard library, run from the directory that contains this note.

\begin{itemize}
\item \texttt{python src/verify.py} checks the Bernoulli polynomials used for \((\chi_3,\chi_3,4)\), the seven quadratics, formula \eqref{eq:eval}, the initial congruences, the power-sum identity of the appendix, and the eighteen words.
The log ends with \texttt{CERTIFICATE OK}.
\item \texttt{python src/check\_words.py certificates/gamma1\_9\_words.txt} and the same program on \texttt{certificates/gamma1\_12\_words.txt} multiply the stored words.
\item \texttt{python src/verify\_chi34.py} checks the nine quadratics, the initial congruences modulo \(3\), and the two span witnesses for level \(12\).
The log ends with \texttt{CHI34 CERTIFICATE OK}.
\item \texttt{python src/verify\_k7.py} checks the degree-\(6\) Bernoulli polynomial, the two level-\(15\) witnesses, the \(192\) cosets, and membership of the \(46\) interpolated polynomials in \(\mathcal{P}(7;4,3)\).
The log ends with \texttt{K7 CERTIFICATE OK}.
\item \texttt{python src/verify\_refinements.py} checks the eighteen triples in Theorem~\ref{thm:mod9}, the \(46\) initial congruences of Section~\ref{sec:15}, and the two quotients of Section~\ref{sec:open}.
It reads the level-\(15\) log, so it is run after \texttt{verify\_k7.py}.
The log ends with \texttt{REFINEMENTS CERTIFICATE OK}.
\item \texttt{python src/verify\_k6.py} checks the weight-\(4\) calibration \(\tS(1,0;9,1)=-2\), the \(72\) cosets and the eighteen words, the eighteen degree-\(4\) polynomials, the input \(\lambda=30\) and \(\alpha=2\), the witnesses \(310/3\) and \(4480/3\), and the parabolic family for \(1\le m\le 6\).
The log ends with \texttt{K6 CERTIFICATE OK}.
\end{itemize}

A finite sample of random words is not one of these inputs and is not a proof.
A greatest common divisor computed on a proper subset of the group is only a lower bound for the span; the spans in Theorem~\ref{thm:three} use a matching containment.

The cocycle proofs and the re-derivation of Lemma~\ref{lem:uni} are Grok's.
The shape of that lemma was first written down by Codex.
The second triple and the value \(-\tfrac{17000}{3}\) were proposed in a separate exploration (Hydrogen) and were recomputed here.
The initial polynomials are taken from \cite{NFDS} and checked against the finite sum wherever a theorem depends on them.
The weight-\(6\) polynomials were recomputed from that finite sum.
The pair \(\lambda=30\), \(\alpha=2\) was first found by Codex.
The arguments were prepared with AI assistance.

\begin{thebibliography}{9}
\bibitem{Majure}
J.~Majure, \emph{Dedekind sums for a Fuchsian group of genus zero}, J.~Number Theory 250 (2023), 35--48.
\bibitem{SVY}
T.~Stucker, A.~Vennos, and M.~P.~Young, \emph{Dedekind sums arising from newform Eisenstein series}, Int.~J.~Number Theory 16 (2020), 2129--2139.
\bibitem{Tran}
P.~Tranbarger, \emph{Higher weight generalized Dedekind sums}, arXiv:2512.17139v1; J.~Number Theory 286 (2026), 1--46.
\bibitem{NFDS}
P.~Tranbarger, repository \texttt{NFDSFastComputation}, \url{https://github.com/prestontranbarger/NFDSFastComputation}.
\end{thebibliography}

\appendix
\section{The values \(-2m^2\)}
\label{app:family}

The span in Section~\ref{sec:9} needs only the case \(m=1\).
The identity for every \(m\) is recorded here because it is exact and short once the block sums are known.

\begin{theorem}
\label{thm:family}
For every integer \(m\ge 1\) and every \(b\in\Z\),
\[
\tS_{\chi_3,\chi_3,4}\begin{pmatrix}1&b\\9m&1+9mb\end{pmatrix}=-2m^2.
\]
\end{theorem}

\begin{proof}
The matrix has determinant \(1\) and lies in \(\Gamma_1(9)\).
For \(c>0\) the finite sum depends only on \(a\) and \(c\), so it is enough to take \(b=0\).
Let \(c=9m\) and \(a=1\).
Terms with \(3\mid j\) vanish.
Write \(j=3q+r\) with \(r\in\{1,2\}\) and \(q\in\{0,\ldots,3m-1\}\), and let \(n\in\{1,2\}\).
Then
\[
x=\frac{j}{c}+\frac{n}{3}=\frac{3(q+nm)+r}{9m}
\]
is never an integer, and \(\lfloor x\rfloor\in\{0,1\}\), with floor \(1\) precisely when \(q\ge m(3-n)\).
There are eight blocks, one for each \((r,n,\lfloor x\rfloor)\).
The summand of \(\tS=c^2 S\) is \(81m^2\,\chi_3(r)\chi_3(n)\,B_1(j/c)\,B_3(x)\).

The substitution \((j,n)\mapsto(c-j,3-n)\) preserves the summand.
Indeed \(c-j=3q'+(3-r)\) with \(q'=3m-1-q\), both character factors change sign because \(\chi_3\) is odd, and \((c-j)/c+(3-n)/3=2-x\).
Thus \(B_1((c-j)/c)=-B_1(j/c)\) and \(B_3(\{2-x\})=-B_3(\{x\})\).
The four signs cancel.
The map identifies the four blocks with \(r=1\) with the four blocks with \(r=2\), and the eight blocks contribute twice the sum of the four blocks with \(r=1\).

On a single block, \(B_3(f)=f(f-\tfrac12)(f-1)\) with \(f=\{x\}\), and the summand equals \(N(q)/(324m^2)\), where \(N\) is the product of the four linear numerators of \(B_1\), \(f\), \(f-\tfrac12\), and \(f-1\), times the character sign.
For the block \((r,n,\lfloor x\rfloor)=(1,1,0)\), summed over \(0\le q\le 2m-1\), the change of coordinate \(t=3q+1\) factors \(N\) as
\[
N=(4t^2-24mt+27m^2)(t^2-3mt-18m^2)
=4t^4-36m\,t^3+27m^2 t^2+351 m^3 t-486 m^4.
\]
In the coordinate \(q\),
\begin{align*}
N={}&324q^4+(432-972m)q^3+(243m^2-972m+216)q^2\\
&+(1053m^3+162m^2-324m+48)q\\
&+(-486m^4+351m^3+27m^2-36m+4).
\end{align*}
Faulhaber's formulae,
\begin{align*}
\sum_{q=0}^{L-1}1&=L,\\
\sum q&=L(L-1)/2,\\
\sum q^2&=L(L-1)(2L-1)/6,\\
\sum q^3&=\bigl(L(L-1)/2\bigr)^2,\\
\sum q^4&=L(L-1)(2L-1)(3L^2-3L-1)/30,
\end{align*}
with \(L=2m\), give \(\sum N=-(162/5)m^5-81m^4+9m^3+24m^2+(52/5)m\).
Dividing by \(324m^2\) produces the block sum
\[
P=-\frac{162m^4+405m^3-45m^2-120m-52}{1620m}.
\]
The numerator factors as \((3m-2)(54m^3+171m^2+99m+26)\), which is checked by expanding the product.
The same power-sum arithmetic on the ranges \(2m\le q\le 3m-1\), \(0\le q\le m-1\), and \(m\le q\le 3m-1\), for the blocks \((1,1,1)\), \((1,2,0)\), and \((1,2,1)\) respectively, produces
\begin{align*}
Q&=\frac{3483m^4-810m^3+315m^2-240m+52}{3240m},\\
R&=-\frac{3483m^4+810m^3+315m^2+240m+52}{3240m},\\
S&=\frac{(3m+2)(54m^3-171m^2+99m-26)}{1620m}.
\end{align*}

The combination is then short.
Write \(X=54m^3+99m\) and \(Y=171m^2+26\).
Then
\[
P+S=\frac{-(3m-2)(X+Y)+(3m+2)(X-Y)}{1620m}=\frac{4X-6mY}{1620m}.
\]
The numerator is \(4X-6mY=216m^3+396m-1026m^3-156m=30m(8-27m^2)\), so
\[
P+S=-\frac{27m^2-8}{54}.
\]
The other pair collapses as
\[
Q+R=\frac{-1620m^3-480m}{3240m}=-\frac{27m^2+8}{54}.
\]
Therefore \(P+Q+R+S=-m^2\), and the eight blocks sum to \(2(P+Q+R+S)=-2m^2\).
\end{proof}

The three closed forms \(Q\), \(R\), and \(S\) are the same kind of finite summation as \(P\).
The program \texttt{src/verify.py} rebuilds each degree-\(4\) numerator from the four linear factors, sums it by the five Faulhaber polynomials, and checks the identity \(2(P+Q+R+S)=-2m^2\) in \(\Q[m]\).
That check agrees with the cancellation above.
As a second check, independent of the block decomposition, the unsplit double sum equals \(-2m^2\) for every \(m=1,\ldots,12\).

\section{Weight \(6\)}
\label{app:k6}

The same group and the same character, at the next even weight, still fit Lemma~\ref{lem:uni}.
The new finite input is one table of eighteen polynomials.
The cocycle induction is not repeated.

\begin{corollary}
\label{cor:k6}
Let \(I\) be the image of \(\tS_{\chi_3,\chi_3,6}\) on \(\Gamma_1(9)\), and set \(U=\tfrac{3}{10}\tS\).
Then \(U(\gamma)\) is an integer and \(U(\gamma)\equiv(c(\gamma)/9)^2\pmod{3}\).
The additive span of \(I\) is \(\tfrac{10}{3}\Z\).
The rationals \(\tfrac{20}{3}\) and \(-\tfrac{10}{3}\) lie in that span and not in \(I\).
The largest positive rational \(r\) with \(I\subseteq r\Z\) is \(\tfrac{10}{3}\).
For every integer \(m\ge 1\) and every \(b\in\Z\),
\begin{equation}
\label{eq:k6family}
\tS_{\chi_3,\chi_3,6}\begin{pmatrix}1&b\\9m&1+9mb\end{pmatrix}
=\frac{10}{3}m^2(27m^2+4).
\end{equation}
\end{corollary}

\begin{proof}
The eighteen non-identity Schreier generators are the stored words for \(\Gamma_1(9)\).
On each of them, \(h=\sum_{i=0}^{4}a_i z^{4-i}\) satisfies \(b_i=9^{i+1}a_i/30\in\Z\) and
\[
b_0\equiv 2r^2,\qquad b_1\equiv -4r,\qquad b_i\equiv 0\quad(i\ge 2)\pmod{3},
\]
with \(r=-c/9\).
Lemma~\ref{lem:uni} applies with \((n,q,\lambda,\alpha)=(4,9,30,2)\).
Here \(-\alpha\equiv 1\pmod{3}\), so \(U(M^{-1})\equiv r(M)^2\pmod{3}\).
For \(\gamma=M^{-1}\) one has \(c(\gamma)/9=r(M)\).
The identity gives the class \(0\).
The matrix \(P=\begin{pmatrix}1&0\\9&1\end{pmatrix}\) has \(\tS(P)=310/3\) and \(\tS(P^2)=4480/3\), hence normalized values \(31\) and \(448\).
These are coprime, so the span is \(\tfrac{10}{3}\Z\), and no larger positive rational has the same containment.
The empty class is \(U\equiv 2\pmod{3}\), which excludes \(\tfrac{20}{3}\) and \(-\tfrac{10}{3}\).
Both surviving classes occur: the identity gives \(0\), and \(U(P)=31\equiv 1\pmod{3}\).

The polynomial of \(P^{-1}=\begin{pmatrix}1&0\\-9&1\end{pmatrix}\) is
\[
h=-\frac{310}{3}x^4+\frac{1160}{27}x^3-\frac{20}{3}x^2+\frac{40}{81}x.
\]
The rescaling \(F(X)=9^5 h(X/9)/30\) is \(-31X^4+116X^3-162X^2+108X\).
Write \(g=P^{-1}\) and \(g_9=\begin{pmatrix}1&0\\-1&1\end{pmatrix}\), so \(g_9^j=\begin{pmatrix}1&0\\-j&1\end{pmatrix}\).
The right cocycle on the rescaled polynomials gives \(F_{g^m}=\sum_{j=0}^{m-1}F|_{-4}g_9^j\).
The matrix \(g^m\) has bottom row \((-9m,1)\), so its ratio is \(r=m\) and its diagonal entry on the right is \(1\).
Therefore \(U(P^m)=m^4 F_{g^m}(1/m)\).
The slash at that point is
\[
\bigl(F|_{-4}g_9^j\bigr)\Bigl(\frac1m\Bigr)=\Bigl(\frac{m-j}{m}\Bigr)^4 F\Bigl(\frac{1}{m-j}\Bigr),
\]
and \(t=m-j\) runs through \(1,\ldots,m\).
Hence
\begin{align*}
U(P^m)
&=\sum_{t=1}^{m}(-31+116t-162t^2+108t^3)\\
&=-31m+58m(m+1)-27m(m+1)(2m+1)+27m^2(m+1)^2\\
&=m^2(27m^2+4).
\end{align*}
Multiplying by \(\tfrac{10}{3}\) gives \eqref{eq:k6family} for \(b=0\).
For \(c>0\) the finite sum depends on the matrix only through its first column, and \(\begin{pmatrix}1&b\\9m&1+9mb\end{pmatrix}\) has the same first column.
\end{proof}

\end{document}
